%!TeX root=../emmatop.tex
\chapter[Chapter \thechapter]{}
\lettrine[lines=4,lraise=0.3]{T}{his} little explanation with Mr Knightley gave Emma considerable pleasure. It was one of the agreeable recollections of the ball, which she walked about the lawn the next morning to enjoy.—She was extremely glad that they had come to so good an understanding respecting the Eltons, and that their opinions of both husband and wife were so much alike; and his praise of Harriet, his concession in her favour, was peculiarly gratifying. The impertinence of the Eltons, which for a few minutes had threatened to ruin the rest of her evening, had been the occasion of some of its highest satisfactions; and she looked forward to another happy result—the cure of Harriet's infatuation.—From Harriet's manner of speaking of the circumstance before they quitted the ballroom, she had strong hopes. It seemed as if her eyes were suddenly opened, and she were enabled to see that Mr Elton was not the superior creature she had believed him. The fever was over, and Emma could harbour little fear of the pulse being quickened again by injurious courtesy. She depended on the evil feelings of the Eltons for supplying all the discipline of pointed neglect that could be farther requisite.—Harriet rational, Frank Churchill not too much in love, and Mr Knightley not wanting to quarrel with her, how very happy a summer must be before her!

She was not to see Frank Churchill this morning. He had told her that he could not allow himself the pleasure of stopping at Hartfield, as he was to be at home by the middle of the day. She did not regret it.

Having arranged all these matters, looked them through, and put them all to rights, she was just turning to the house with spirits freshened up for the demands of the two little boys, as well as of their grandpapa, when the great iron sweep-gate opened, and two persons entered whom she had never less expected to see together—Frank Churchill, with Harriet leaning on his arm—actually Harriet!—A moment sufficed to convince her that something extraordinary had happened. Harriet looked white and frightened, and he was trying to cheer her.—The iron gates and the front-door were not twenty yards asunder;—they were all three soon in the hall, and Harriet immediately sinking into a chair fainted away.

A young lady who faints, must be recovered; questions must be answered, and surprizes be explained. Such events are very interesting, but the suspense of them cannot last long. A few minutes made Emma acquainted with the whole.

Miss Smith, and Miss Bickerton, another parlour boarder at Mrs Goddard's, who had been also at the ball, had walked out together, and taken a road, the Richmond road, which, though apparently public enough for safety, had led them into alarm.—About half a mile beyond Highbury, making a sudden turn, and deeply shaded by elms on each side, it became for a considerable stretch very retired; and when the young ladies had advanced some way into it, they had suddenly perceived at a small distance before them, on a broader patch of greensward by the side, a party of gipsies. A child on the watch, came towards them to beg; and Miss Bickerton, excessively frightened, gave a great scream, and calling on Harriet to follow her, ran up a steep bank, cleared a slight hedge at the top, and made the best of her way by a short cut back to Highbury. But poor Harriet could not follow. She had suffered very much from cramp after dancing, and her first attempt to mount the bank brought on such a return of it as made her absolutely powerless—and in this state, and exceedingly terrified, she had been obliged to remain.

\begin{figure}[tbph]
\centering
\includegraphics[width=.8\linewidth]{39assailed}
\caption{Harriet was soon assailed}
\end{figure}

How the trampers might have behaved, had the young ladies been more courageous, must be doubtful; but such an invitation for attack could not be resisted; and Harriet was soon assailed by half a dozen children, headed by a stout woman and a great boy, all clamorous, and impertinent in look, though not absolutely in word.—More and more frightened, she immediately promised them money, and taking out her purse, gave them a shilling, and begged them not to want more, or to use her ill.—She was then able to walk, though but slowly, and was moving away—but her terror and her purse were too tempting, and she was followed, or rather surrounded, by the whole gang, demanding more.

In this state Frank Churchill had found her, she trembling and conditioning, they loud and insolent. By a most fortunate chance his leaving Highbury had been delayed so as to bring him to her assistance at this critical moment. The pleasantness of the morning had induced him to walk forward, and leave his horses to meet him by another road, a mile or two beyond Highbury—and happening to have borrowed a pair of scissors the night before of Miss Bates, and to have forgotten to restore them, he had been obliged to stop at her door, and go in for a few minutes: he was therefore later than he had intended; and being on foot, was unseen by the whole party till almost close to them. The terror which the woman and boy had been creating in Harriet was then their own portion. He had left them completely frightened; and Harriet eagerly clinging to him, and hardly able to speak, had just strength enough to reach Hartfield, before her spirits were quite overcome. It was his idea to bring her to Hartfield: he had thought of no other place.

This was the amount of the whole story,—of his communication and of Harriet's as soon as she had recovered her senses and speech.—He dared not stay longer than to see her well; these several delays left him not another minute to lose; and Emma engaging to give assurance of her safety to Mrs Goddard, and notice of there being such a set of people in the neighbourhood to Mr Knightley, he set off, with all the grateful blessings that she could utter for her friend and herself.

Such an adventure as this,—a fine young man and a lovely young woman thrown together in such a way, could hardly fail of suggesting certain ideas to the coldest heart and the steadiest brain. So Emma thought, at least. Could a linguist, could a grammarian, could even a mathematician have seen what she did, have witnessed their appearance together, and heard their history of it, without feeling that circumstances had been at work to make them peculiarly interesting to each other?—How much more must an imaginist, like herself, be on fire with speculation and foresight!—especially with such a groundwork of anticipation as her mind had already made.

It was a very extraordinary thing! Nothing of the sort had ever occurred before to any young ladies in the place, within her memory; no rencontre, no alarm of the kind;—and now it had happened to the very person, and at the very hour, when the other very person was chancing to pass by to rescue her!—It certainly was very extraordinary!—And knowing, as she did, the favourable state of mind of each at this period, it struck her the more. He was wishing to get the better of his attachment to herself, she just recovering from her mania for Mr Elton. It seemed as if every thing united to promise the most interesting consequences. It was not possible that the occurrence should not be strongly recommending each to the other.

In the few minutes' conversation which she had yet had with him, while Harriet had been partially insensible, he had spoken of her terror, her naïveté, her fervour as she seized and clung to his arm, with a sensibility amused and delighted; and just at last, after Harriet's own account had been given, he had expressed his indignation at the abominable folly of Miss Bickerton in the warmest terms. Every thing was to take its natural course, however, neither impelled nor assisted. She would not stir a step, nor drop a hint. No, she had had enough of interference. There could be no harm in a scheme, a mere passive scheme. It was no more than a wish. Beyond it she would on no account proceed.

Emma's first resolution was to keep her father from the knowledge of what had passed,—aware of the anxiety and alarm it would occasion: but she soon felt that concealment must be impossible. Within half an hour it was known all over Highbury. It was the very event to engage those who talk most, the young and the low; and all the youth and servants in the place were soon in the happiness of frightful news. The last night's ball seemed lost in the gipsies. Poor Mr Woodhouse trembled as he sat, and, as Emma had foreseen, would scarcely be satisfied without their promising never to go beyond the shrubbery again. It was some comfort to him that many inquiries after himself and Miss Woodhouse (for his neighbours knew that he loved to be inquired after), as well as Miss Smith, were coming in during the rest of the day; and he had the pleasure of returning for answer, that they were all very indifferent—which, though not exactly true, for she was perfectly well, and Harriet not much otherwise, Emma would not interfere with. She had an unhappy state of health in general for the child of such a man, for she hardly knew what indisposition was; and if he did not invent illnesses for her, she could make no figure in a message.

The gipsies did not wait for the operations of justice; they took themselves off in a hurry. The young ladies of Highbury might have walked again in safety before their panic began, and the whole history dwindled soon into a matter of little importance but to Emma and her nephews:—in her imagination it maintained its ground, and Henry and John were still asking every day for the story of Harriet and the gipsies, and still tenaciously setting her right if she varied in the slightest particular from the original recital.